\documentclass[10pt,a4paper]{amsart}
\usepackage{amsmath,amssymb,amsthm,graphicx}
% Retained preamble for the Stacks Project chapter {CH} (shared project preamble).
% The chapter file pinned-original-source.tex only carries \input{{preamble}} on its first line,
% so no preamble text lives inside the chapter itself.  The Stacks Project uses one single shared
% preamble for every chapter: the snapshot stages it three times (sources/stacks-chow/preamble.tex,
% sources/stacks-etale-cohomology/preamble.tex and focused-seeds/AGP/sources/stacks-snapshot-20261011/
% preamble.tex) and those copies are identical in size.  The block below is copied mechanically,
% line for line, from sources/stacks-chow/preamble.tex lines 14-18 and 154-217: the Xy-pic setup and
% the native theorem declarations together with the mathematical macros.  Nothing was added, removed
% or redefined; the wrapper class and the amsmath/amssymb/amsthm/graphicx base packages are supplied
% by the assembler.
% For commutative diagrams we use Xy-pic
% Chapter: stacks-cohomology.tex -- Cohomology of Sheaves.
\usepackage[all]{xy}

% We use 2cell for 2-commutative diagrams.
\xyoption{2cell}

% Theorem environments.
%
\theoremstyle{plain}
\newtheorem{theorem}[subsection]{Theorem}
\newtheorem{proposition}[subsection]{Proposition}
\newtheorem{lemma}[subsection]{Lemma}

\theoremstyle{definition}
\newtheorem{definition}[subsection]{Definition}
\newtheorem{example}[subsection]{Example}
\newtheorem{exercise}[subsection]{Exercise}
\newtheorem{situation}[subsection]{Situation}

\theoremstyle{remark}
\newtheorem{remark}[subsection]{Remark}
\newtheorem{remarks}[subsection]{Remarks}

\numberwithin{equation}{subsection}

% Macros
%
\def\lim{\mathop{\mathrm{lim}}\nolimits}
\def\colim{\mathop{\mathrm{colim}}\nolimits}
\def\Spec{\mathop{\mathrm{Spec}}}
\def\Hom{\mathop{\mathrm{Hom}}\nolimits}
\def\Ext{\mathop{\mathrm{Ext}}\nolimits}
\def\SheafHom{\mathop{\mathcal{H}\!\mathit{om}}\nolimits}
\def\SheafExt{\mathop{\mathcal{E}\!\mathit{xt}}\nolimits}
\def\Sch{\mathit{Sch}}
\def\Mor{\mathop{\mathrm{Mor}}\nolimits}
\def\Ob{\mathop{\mathrm{Ob}}\nolimits}
\def\Sh{\mathop{\mathit{Sh}}\nolimits}
\def\NL{\mathop{N\!L}\nolimits}
\def\CH{\mathop{\mathrm{CH}}\nolimits}
\def\proetale{{pro\text{-}\acute{e}tale}}
\def\etale{{\acute{e}tale}}
\def\QCoh{\mathit{QCoh}}
\def\Ker{\mathop{\mathrm{Ker}}}
\def\Im{\mathop{\mathrm{Im}}}
\def\Coker{\mathop{\mathrm{Coker}}}
\def\Coim{\mathop{\mathrm{Coim}}}

% Boxtimes
%
\DeclareMathSymbol{\boxtimes}{\mathbin}{AMSa}{"02}

%
% Macros for moduli stacks/spaces
%
\def\QCohstack{\mathcal{QC}\!\mathit{oh}}
\def\Cohstack{\mathcal{C}\!\mathit{oh}}
\def\Spacesstack{\mathcal{S}\!\mathit{paces}}
\def\Quotfunctor{\mathrm{Quot}}
\def\Hilbfunctor{\mathrm{Hilb}}
\def\Curvesstack{\mathcal{C}\!\mathit{urves}}
\def\Polarizedstack{\mathcal{P}\!\mathit{olarized}}
\def\Complexesstack{\mathcal{C}\!\mathit{omplexes}}
% \Pic is the operator that assigns to X its picard group, usage \Pic(X)
% \Picardstack_{X/B} denotes the Picard stack of X over B
% \Picardfunctor_{X/B} denotes the Picard functor of X over B
\def\Pic{\mathop{\mathrm{Pic}}\nolimits}
\def\Picardstack{\mathcal{P}\!\mathit{ic}}
\def\Picardfunctor{\mathrm{Pic}}
\def\Deformationcategory{\mathcal{D}\!\mathit{ef}}

% --- appended 2026-10-12 (Chief): the project own definitions of the reference/slogan/history
% display environments, copied mechanically from the shared Stacks preamble lines 8-12
% (\usepackage{verbatim} plus the three \newenvironment lines, which the chapter preamble keeps
% verbatim from the project preamble). Nothing added, removed or redefined: these environments
% carry no mathematical content and the project itself typesets them as comments. Without these
% three lines a counted statement whose own text uses them fails with Environment slogan undefined.
\usepackage{verbatim}
\newenvironment{reference}{\comment}{\endcomment}
\newenvironment{slogan}{\comment}{\endcomment}
\newenvironment{history}{\comment}{\endcomment}
\usepackage{cleveref}
\crefname{originalremark}{original remark}{original remarks}
\Crefname{originalremark}{Original remark}{Original remarks}
\makeatletter
\newlabel{remark-prepare-fancy-base-change}{{08I3}{source}{}{sourcepointer-1}{}}
\newlabel{remark-prepare-fancy-base-change@cref}{{[originalremark][0][]08I3}{[1][0][]source}{}{}{}}
\makeatother
\crefname{originalsection}{original section}{original sections}
\Crefname{originalsection}{Original section}{Original sections}
\makeatletter
\newlabel{section-perfect}{{08CL}{source}{}{sourcepointer-2}{}}
\newlabel{section-perfect@cref}{{[originalsection][0][]08CL}{[1][0][]source}{}{}{}}
\makeatother
\crefname{originallemma}{original lemma}{original lemmas}
\Crefname{originallemma}{Original lemma}{Original lemmas}
\makeatletter
\newlabel{lemma-Rhom-strictly-perfect}{{08DM}{source}{}{sourcepointer-3}{}}
\newlabel{lemma-Rhom-strictly-perfect@cref}{{[originallemma][0][]08DM}{[1][0][]source}{}{}{}}
\makeatother
\crefname{originalsection}{original section}{original sections}
\Crefname{originalsection}{Original section}{Original sections}
\makeatletter
\newlabel{section-derived-pullback}{{06YI}{source}{}{sourcepointer-4}{}}
\newlabel{section-derived-pullback@cref}{{[originalsection][0][]06YI}{[1][0][]source}{}{}{}}
\makeatother
\crefname{originallemma}{original lemma}{original lemmas}
\Crefname{originallemma}{Original lemma}{Original lemmas}
\makeatletter
\newlabel{lemma-strictly-perfect-pullback}{{09U6}{source}{}{sourcepointer-5}{}}
\newlabel{lemma-strictly-perfect-pullback@cref}{{[originallemma][0][]09U6}{[1][0][]source}{}{}{}}
\makeatother
\crefname{originallemma}{original lemma}{original lemmas}
\Crefname{originallemma}{Original lemma}{Original lemmas}
\makeatletter
\newlabel{lemma-Rhom-strictly-perfect-K-flat}{{0GM5}{source}{}{sourcepointer-6}{}}
\newlabel{lemma-Rhom-strictly-perfect-K-flat@cref}{{[originallemma][0][]0GM5}{[1][0][]source}{}{}{}}
\makeatother
\crefname{originalsection}{original section}{original sections}
\Crefname{originalsection}{Original section}{Original sections}
\makeatletter
\newlabel{section-pseudo-coherent}{{08CA}{source}{}{sourcepointer-7}{}}
\newlabel{section-pseudo-coherent@cref}{{[originalsection][0][]08CA}{[1][0][]source}{}{}{}}
\makeatother
\title{Base change for the internal hom: pullback of R-Hom is R-Hom of the pullbacks in three cases}
\author{The Stacks Project Authors}
\date{stacks-project:cohomology.tex:local-snapshot-20261011; GNU FDL 1.2 or later}
\begin{document}
\maketitle
\noindent\textit{Editorial note.} All delivered mathematical text is the author's own, copied verbatim from the pinned Stacks Project chapter file cohomology.tex; no statement or proof step was rewritten, reordered or removed.\par\medskip
\noindent\textit{Editorial note.} The counted span cites seven cross-chapter labels (including a remark and a section), each recorded in assembly.external\_references with its canonical tag as a navigation pointer.\par\medskip
\noindent\textit{Editorial note.} The proof cites no bibliography item inside the counted span.\par\medskip
\noindent\textit{Allowed standard prerequisites (editorial).}\begin{itemize}
\item Computation of R-Hom out of a strictly perfect complex. Stacks Project, Cohomology of Sheaves, Lemma lemma-Rhom-strictly-perfect (tag 0A8D). Represents the internal hom by an explicit complex of Hom sheaves.
\item Pullback of finite free and strictly perfect complexes. Stacks Project, Cohomology of Sheaves, Lemma lemma-strictly-perfect-pullback (tag 09U6). Shows the comparison map is computed levelwise on free modules.
\end{itemize}\medskip
\ifdefined\hypertarget\hypertarget{sourcepointer-1}{}\fi
\noindent\textit{Source reference.} Original remark 08I3: sources/stacks-chow/tags.txt, source lines 10801--10801. Cross-chapter standard prerequisite cited by the author inside the counted span; it is used as a pointer and its own proof is not part of this item.\par\medskip
\ifdefined\hypertarget\hypertarget{sourcepointer-2}{}\fi
\noindent\textit{Source reference.} Original section 08CL: sources/stacks-chow/tags.txt, source lines 10606--10606. Cross-chapter standard prerequisite cited by the author inside the counted span; it is used as a pointer and its own proof is not part of this item.\par\medskip
\ifdefined\hypertarget\hypertarget{sourcepointer-3}{}\fi
\noindent\textit{Source reference.} Original lemma 08DM: sources/stacks-chow/tags.txt, source lines 10642--10642. Cross-chapter standard prerequisite cited by the author inside the counted span; it is used as a pointer and its own proof is not part of this item.\par\medskip
\ifdefined\hypertarget\hypertarget{sourcepointer-4}{}\fi
\noindent\textit{Source reference.} Original section 06YI: sources/stacks-chow/tags.txt, source lines 8887--8887. Cross-chapter standard prerequisite cited by the author inside the counted span; it is used as a pointer and its own proof is not part of this item.\par\medskip
\ifdefined\hypertarget\hypertarget{sourcepointer-5}{}\fi
\noindent\textit{Source reference.} Original lemma 09U6: sources/stacks-chow/tags.txt, source lines 12414--12414. Cross-chapter standard prerequisite cited by the author inside the counted span; it is used as a pointer and its own proof is not part of this item.\par\medskip
\ifdefined\hypertarget\hypertarget{sourcepointer-6}{}\fi
\noindent\textit{Source reference.} Original lemma 0GM5: sources/stacks-chow/tags.txt, source lines 20745--20745. Cross-chapter standard prerequisite cited by the author inside the counted span; it is used as a pointer and its own proof is not part of this item.\par\medskip
\ifdefined\hypertarget\hypertarget{sourcepointer-7}{}\fi
\noindent\textit{Source reference.} Original section 08CA: sources/stacks-chow/tags.txt, source lines 10595--10595. Cross-chapter standard prerequisite cited by the author inside the counted span; it is used as a pointer and its own proof is not part of this item.\par\medskip
\begin{lemma}
\label{lemma-pullback-internal-hom}
Let $h : X \to Y$ be a morphism of ringed spaces.
Let $K, M$ be objects of $D(\mathcal{O}_Y)$. The
canonical map
$$
Lh^*R\SheafHom(K, M) \longrightarrow R\SheafHom(Lh^*K, Lh^*M)
$$
of Remark \ref{remark-prepare-fancy-base-change}
is an isomorphism in the following cases
\begin{enumerate}
\item $K$ is perfect,
\item $h$ is flat, $K$ is pseudo-coherent, and $M$ is (locally) bounded below,
\item $\mathcal{O}_X$ has finite tor dimension over $h^{-1}\mathcal{O}_Y$,
$K$ is pseudo-coherent, and $M$ is (locally) bounded below,
\end{enumerate}
\end{lemma}

\begin{proof}
Proof of (1). The question is local on $Y$, hence we may assume that
$K$ is represented by a strictly perfect complex
$\mathcal{E}^\bullet$, see Section \ref{section-perfect}.
Choose a K-flat complex $\mathcal{F}^\bullet$ representing $M$.
Apply Lemma \ref{lemma-Rhom-strictly-perfect} to see that
$R\SheafHom(K, L)$ is represented by the complex
$\mathcal{H}^\bullet =
\SheafHom^\bullet(\mathcal{E}^\bullet, \mathcal{F}^\bullet)$
with terms $\mathcal{H}^n = \bigoplus\nolimits_{n = p + q}
\SheafHom_{\mathcal{O}_X}(\mathcal{E}^{-q}, \mathcal{F}^p)$.
By the construction of $Lh^*$ in Section \ref{section-derived-pullback}
we see that $Lh^*K$ is represented by the strictly perfect complex
$h^*\mathcal{E}^\bullet$ (Lemma \ref{lemma-strictly-perfect-pullback}).
Similarly, the object $Lh^*M$ is represented by
the complex $h^*\mathcal{F}^\bullet$.
Finally, the object $Lh^*R\SheafHom(K, M)$
is represented by $h^*\mathcal{H}^\bullet$ as
$\mathcal{H}^\bullet$ is K-flat by
Lemma \ref{lemma-Rhom-strictly-perfect-K-flat}.
Thus to finish the proof it suffices to show that
$h^*\mathcal{H}^\bullet =
\SheafHom^\bullet(h^*\mathcal{E}^\bullet, h^*\mathcal{F}^\bullet)$.
For this it suffices to note that
$h^*\SheafHom(\mathcal{E}, \mathcal{F}) =
\SheafHom(h^*\mathcal{E}, \mathcal{F})$
whenever $\mathcal{E}$ is a direct summand of a finite free
$\mathcal{O}_X$-module.

\medskip\noindent
Proof of (2). Since $h$ is flat, we can compute $Lh^*$
by simply using $h^*$ on any complex of $\mathcal{O}_Y$-modules.
In particular we have $H^i(Lh^*K) = h^*H^i(K)$ for all $i \in \mathbf{Z}$.
Say $H^i(M) = 0$ for $i < a$. Let $K' \to K$ be a morphism
of $D(\mathcal{O}_Y)$ which defines an isomorphism
$H^i(K') \to H^i(K)$ for all $i \geq b$. Then the corresponding maps
$$
R\SheafHom(K, M) \to R\SheafHom(K', M)
$$
and
$$
R\SheafHom(Lh^*K, Lh^*M) \to R\SheafHom(Lh^*K', Lh^*M)
$$
are isomorphisms on cohomology sheaves in degrees $< a - b$ (details omitted).
Thus to prove the map in the statement of the lemma induces an
isomorphism on cohomology sheaves in degrees $< a - b$ it suffices to
prove the result for $K'$ in those degrees. Also, as in the proof
of part (1) the question is local on $Y$. Thus we may assume $K$ is
represented by a strictly perfect complex, see
Section \ref{section-pseudo-coherent}. This reduces us to case (1).

\medskip\noindent
Proof of (3). The proof is the same as the proof of (2) except one
uses that $Lh^*$ has bounded cohomological dimension to get the
desired vanishing. We omit the details.
\end{proof}

\end{document}
